\documentclass[10pt,oneside,headheight=10mm]{article}
% Packages
\usepackage{amsmath}    % American Math Socitey distributed package
\usepackage{amsfonts}	% Ams Fonts
\usepackage{mathrsfs}	% Ams Fonts
\usepackage{amssymb}	% Ams Symbol
\usepackage[a4paper, left=1.2cm,right=1.5cm,top=2cm,bottom=2cm]{geometry}   % Page size and margin
\usepackage{ulem}   % Package for underlines
\usepackage{graphicx}   % Package for including images
\usepackage{color}  % Font color
\usepackage[x11names]{xcolor}   % Adds more color and enables customize colors
\usepackage{afterpage}   % Set multiple pages in one line
\usepackage{utfsym} % More symbols
\usepackage{titlesec}   % Program title format for environments
\usepackage{sectsty}    % Section title color
\usepackage{fancyhdr}   % Make headers and footers
\usepackage{tipa}	%ipa chart
\usepackage{mathtools}	% More Symbols
\usepackage{hyperref} 	% For links, urls
\usepackage{listings}	% Including codes
\usepackage{multicol}
\usepackage{pdfpages}


\usepackage{pgfplots}
\pgfplotsset{width=8cm,compat=1.9}


% Other Settings
\setlength{\fboxrule}{1.2pt}    % Setting thickness of fbox
\graphicspath{ {./images/} }    % Setting defulat image path
\definecolor{umYellow}{HTML}{FFCB05}    % Defining color - umYellow
\definecolor{umBlue}{HTML}{00274C}      % Defining color - umBlue

% Set Environment Appearance & New enviornments
\allsectionsfont{\color{umBlue}}
\newtheorem{theorem}{\color{umBlue}Theorem}[section]
\newtheorem{example}{\color{umBlue}Example}[section]
\newtheorem{definition}{\color{umBlue}Definition}[section]
\newtheorem{corollary}{\color{umBlue}Corollary}[theorem]
\newtheorem{lemma}{\color{umBlue}Lemma}[theorem]

\setlength{\parskip}{0.254cm} % Distance Between Paragraphs
\setlength{\parindent}{0pt}    % Indent first line of parargraph



% New Commands
\newcommand{\bm}[1]{\begin{bmatrix}#1\end{bmatrix}}
\newcommand{\cb}[1]{\left(#1\right)}
\newcommand{\sqb}[1]{\left[#1\right]}
\newcommand{\pp}{\partial}
\newcommand{\tx}[1]{\text{#1}}
\newcommand{\vs}[1]{\vspace{#1}}
\newcommand{\solution}{\color{umBlue}$~~~~$\textit{solution.}\color{black}\\}
\newcommand{\proof}{\color{umBlue}$~~~~$\textit{proof.}\color{black}\\}
\newcommand{\qed}{\hfill$\blacksquare$}
\newcommand{\cut}{\backslash}
\newcommand{\bra}[1]{\left<#1\right|}
\newcommand{\ket}[1]{\left|#1\right>}
\newcommand{\hs}{\hspace*{1cm}}
\newcommand{\ib}[1]{\left.#1\right|}
\newcommand{\dbar}{\text{\dj}}
\usepackage{mahjong}

\numberwithin{equation}{section}

% Document

\begin{document}

\begin{center}
\fcolorbox{umBlue}{umBlue}{\parbox{\linewidth}{
\begin{center}
\LARGE \color{Chocolate1}\textbf{Numerical Solution to the Perturbation Equations}

\vspace{-0.2cm}
{\color{Chocolate1}\rule{14cm}{0.6pt}}\\\color{white}

\vspace{0.2cm}
\large chih-uva\\
\today

\end{center}}}
\end{center}
The purpose of writing this document is to stepwisely derive the numerical solution of the set of perturbation equations of the TLN-TDN project.


\section{Solution to the TOV Equation}
TOV equation is a set of equation that describes the internal structure of the star, given a equation of state (EOS), which is a function that describes the relation between stellar pressure and density. 

The TOV equation has the form:
\begin{align}
    \frac{dp}{dr} = -\frac{\cb{p+\rho}\cb{m+4\pi r^3 p}}{r(r-2m)}
\end{align}

Given a background metric:
\begin{align}
    ds^2 = -e^\nu dt^2 + e^{\lambda} dr^2 + r^2d\Omega^2,\nonumber
\end{align}
one is also able to obtain the relation:
\begin{align}
    \frac{d\nu}{dr} = 2\frac{m+ 4\pi r^3 p}{r(r-2m)}
\end{align}
Also note that we have the conservation equation:
\begin{align}
    \frac{dm}{dr} = 4\pi r^2 \rho
\end{align}
With this set of equation, we are able to solve the background stellar pressure and mass functions. This is a set of first-order ODEs, however, their initial condition are not independent.

We assume functional forms of the TOV equations to be polynomials, with max order 3 (by omitting the higher-order contribution, and the $r$ contribution to the mass function is in the third order):
\begin{align}
    \nu(r) =&\, a_0 + a_1r + a_2r^2 + a_3r^3\\
    m(r) =&\, b_1r + b_2r^2 + b_3r^3\\
    p(r) =&\, c_0 + c_1r + c_2r^2 + c_3r^3
\end{align}
We have set $b_0 = 0$ as the initial condition for the mass function.

Substitute the polynomials back to the set of equations, and expand them to the second order of $r$. The mass conservation yields:
\begin{align}
    b_1+r \left(2 b_2 + 3b_3 r-4 \pi  r \left(\frac{c_0}{K}\right)^{n/(n+1)}\right)=0,
\end{align}
Which indicates:
\begin{align}\label{eqn:1-7}
    b_1 = b_2 = 0, b_3 = \frac{4\pi}{3}\cb{\frac{c_0}{K}}^{n/(n+1)}
\end{align}
Solve for $c_0$ we obtain:
\begin{align}
    c_0 = K\cb{\frac{3b_3}{4\pi}}^{1+\frac{1}{n}}
\end{align}
Given that $b_1 = b_2 = 0$, the $d\nu/dr$ relation gives:
\begin{align}
    a_1+r (2 a_2-2 b_3-8 \pi  c_0)+r^2 (3 a_3-8 \pi  c_1) = 0
\end{align}
We have:
\begin{align}
    a_1 = 0, \,\,c_0 = \frac{a_2 - b_2}{4\pi},\,\,c_1 = \frac{3}{8\pi}a_3
\end{align}
With $a_1 = 0$, we obtain:
\begin{align}\label{eqn:1-12}
  & c_1 + r \left(
    b_3 \left(\frac{c_0}{K}\right)^{\frac{n}{n+1}}
    + b_3 c_0
    + 4\pi c_0^2
    + 4\pi c_0 \left(\frac{c_0}{K}\right)^{\frac{n}{n+1}}
    + 2c_2
\right)\nonumber\\
 & + r^2 \left(
    \frac{b_3 c_1 n \left(\frac{c_0}{K}\right)^{\frac{n}{n+1}}}{c_0 (n+1)}
    + b_3 c_1
    + \frac{4\pi c_1 (2n+1) \left(\frac{c_0}{K}\right)^{\frac{n}{n+1}}}{n+1}
    + 8\pi c_0 c_1
    + 3c_3
\right)= 0
\end{align}
We observe that $c1 = 0$ and therefore $a_3 = 0$. We substitute this result back into \autoref{eqn:1-12} and obtain:
\begin{align}
    r\, \cb{ b_3 \left( \frac{c_0}{K} \right)^{\frac{n}{n+1}} + b_3 c_0 + 4\pi c_0^2 + 4\pi c_0 \left( \frac{c_0}{K} \right)^{\frac{n}{n+1}} + 2c_2 } + 3c_3 r^2 = 0
\end{align}
We obtain that $c_3 = 0$ and a rather complicated relation between $c_0, c_2, b_3$, but with \autoref{eqn:1-7}:
\begin{align}
    c_2 = -\frac{1}{32} \pi c_0 \left(9 \pi c_0 K + 64\right) \cb{ \left( \frac{c_0}{K} \right)^{\frac{n}{n+1}} + c_0 }
\end{align}
To summarize:
\begin{align}
    b_1 =&\,\, b_2 = c_1 = 0;\hspace{0.3cm} c_0 = K\cb{\frac{3b_3}{4\pi}}^{1+\frac{1}{n}}  =\frac{a_2 - b_2}{4\pi}\\
    c_2 =& -\frac{1}{32} \pi c_0 \left(9 \pi c_0 K + 64\right) \cb{ \left( \frac{c_0}{K} \right)^{\frac{n}{n+1}} + c_0 }
\end{align}
With the hidden relation:
\begin{align}
    \frac{dp}{dr} = -\frac{1}{2\cb{p+\rho}}\frac{d\nu}{dr},
\end{align}
we are also able to obtain:
\begin{align}
    a_1 = 0, \hspace{0.4cm}a_2 + 2c_2 \cb{ \left( \frac{c_0}{K} \right)^{\frac{n}{n+1}} + c_0 } = 0, \hspace{0.4cm} a_3 + 2c_3 \left( c_0 + \left( \frac{c_0}{K} \right)^{\frac{n}{n+1}} \right) = 0
\end{align}
We are then able to reduce the ansatz to the form:
\begin{align}
    \nu(r) =&\, a_0 + a_2r^2 + a_3r^3\\
    m(r) =&\,  \frac{4\pi}{3}\cb{\frac{c_0}{K}}^{n/(n+1)}r^3\\
    p(r) =&\, c_0 + c_2r^2 + c_3r^3
\end{align}
\rule{\linewidth}{0.8pt}
I think it's not really possible to obtain a relation between $a_0$ and $c_0$ from this set of equations, as there is no $\nu$ term and $a_0$ is always eliminated in the $\nu'$ term.


To match this result,  we compare to the asymptotic behaviors in \href{https://arxiv.org/pdf/1302.1918}{\color{Chocolate1} 1302.1918}.
\begin{align}
    \rho(R) =& \rho_c + \rho_2 R^2 + \mathcal{O}(R^3)\\
    p(R) =& p_c + \frac{2\pi}{3}(\rho_c + p_c)(\rho_c + 3p_c)R^2 + \mathcal{O}\cb{R^3}\\
    M(R) =& \frac{4\pi}{3}R^3 + \frac{4\pi}{5}\rho_2 R^5 + \mathcal{O}\cb{R^6}\\
    \nu(R) =& \nu_c + \frac{4\pi}{3}\cb{\rho_c + 3p_c}R^2 + \mathcal{O}(R^3)
\end{align}

We recognize that with the polytropic EOS, $\frac{4\pi}{3}\cb{p_0/K}^{n/(n+1)}$ is the central density, and the terms of my calculation in general matches the result in the paper. 






\subsection{A dimensionless equation?}
Since there aren't fixed relation between $a_0$ and $c_0$, but we still want initial condition for $\nu$. We therefore consider the stellar surface, where the exterior solution satisfies the birkhoff's theorem, meaning the external metric satisfies the Schwarzschild solution.

We introduce new variables:
\begin{align}
    \bar{r} :=\frac{r}{R}, \hspace{1cm}\bar{m} :=\frac{m}{M},\hspace*{1cm}\bar{P} = \frac{p}{p_c}
\end{align}
With a polytropic EOS:
\begin{align}
    p = K\rho^{1+1/n} \hspace*{1cm} \Rightarrow\hspace*{1cm} \rho = \cb{\frac{p}{K}}^{n/(n+1)}
\end{align}



\section{Solution to the set of perturbation equations}
The set of perturbation so solve attains the form:
\begin{align}
H_1'=&\frac{e^\lambda}{r} H+\left[4\pi r \left(\rho-p\right)e^\lambda -\frac{2m e^\lambda}{r^2}-\frac{\ell+1}{r}\right]H_1+\frac{e^\lambda}{r} K-16 \pi\frac{\rho+p}{r}e^\lambda V-32\pi i  \omega \eta \frac{e^{\lambda-\nu/2}}{r}V,\label{eq:eqforH1}\\
K'=&\frac{H}{r}+\frac{\ell\left(\ell+1\right)}{2r}H_1+\left(\frac{\nu'}{2}-\frac{\ell+1}{r}\right)K-8\pi\frac{\rho+p}{r}e^{\lambda/2} W-32\pi i\omega \eta \frac{e^{-\nu/2}}{r}V,\label{eq:eqforK}\\
W'=&\frac{r}{2}e^{\lambda/2}H+r e^{\lambda/2}K+\frac{r e^{(\lambda-\nu)/2}}{\rho+p}c_s^{-2}X-\frac{\ell+1}{r}W -\frac{\ell\left(\ell+1\right)}{r}e^{\lambda/2}V-16 \pi i \omega \eta r e^{(\lambda-\nu)/2}V,\label{eq:eqforW}\\
X'=&\frac{\rho+p}{2}e^{\nu/2}\left(\frac{1}{r}-\frac{\nu'}{2}\right)H+\frac{\rho+p}{2}e^{\nu/2}\left[\omega^2 r e^{-\nu}+\frac{\ell\left(\ell+1\right)}{2r}\right]H_1+\frac{\rho+p}{2}e^{\nu/2}\left(\frac{3\nu'}{2}-\frac{1}{r}\right)K  \nonumber\\
&-\frac{\ell}{r}X-\frac{\rho+p}{r}e^{(\lambda+\nu)/2}\left[\omega^2 e^{-\nu}+\frac{e^{-\lambda}}{4r^2}\left\{e^{2\lambda}\left(1+8\pi r^2 p \right)^2+2e^\lambda\left(3+8\pi r^2 p \right)-7\right\}\right]W+\frac{\ell\left(\ell+1\right)}{r^2}e^{\nu/2}p' V \label{eq:eqforpX}\\
&+\alpha_\eta \eta+\alpha_{\eta'} \eta'+\alpha_\zeta \zeta+\alpha_{\zeta'} \zeta'\nonumber
\end{align}

Where we have above solved the numerical profile for $\rho, p, \nu$ and $\lambda$.








\end{document}
